% --- Archon physics formalization source begin ---
% archon:physics
% archon:covers IPhO2026Problems/problem_IPhO_2026_1_C_2.lean
% archon:source-report reports/ipho_2026_k3/problem_IPhO_2026_1_C_2.source.json
% archon:problem-id IPhO_2026_1
% archon:part-id C.2
% archon:previous-part IPhO_2026_1_C_1
% archon:previous-part-policy natural_language_prerequisite_only

\chapter{Physics problem IPhO\_2026\_1\_C\_2}
\label{ch:IPhO2026Problems_problem_IPhO_2026_1_C_2}

\paragraph{Problem source.}
A photon of angular frequency omega is absorbed by an ozone molecule O3 at rest,
dissociating it into O2 and O.  Let U\_i and U\_f be the ground-state energies of
O3 and O2 and define Delta U = U\_f - U\_i.  The outgoing O2 momentum makes angle
theta with the incident photon.  Treat the oxygen fragments classically and
non-relativistically, take the mass of an oxygen atom to be m, and use photon
momentum p\_gamma = E\_gamma/c = hbar*omega/c.

Current subquestion:
For theta = pi/6, Delta U = 1.10 eV, and m = 16.0 amu, calculate hbar*omega\_min - Delta U in eV.

\paragraph{Current subquestion.}
For theta = pi/6, Delta U = 1.10 eV, and m = 16.0 amu, calculate hbar*omega\_min - Delta U in eV.

\paragraph{Recorded answer/context.}
hbar*omega\_min - Delta U = 2.03e-11 eV.

\paragraph{Figure/image path.}
/root/proposal\_for\_physic/science-mango/ipho\_2026\_source/image/T1\_page-3.png

\paragraph{Reusable previous-part conclusions.}
\begin{itemize}
\item Source C.1. Question: Determine the minimum angular frequency omega\_min required for dissociation at outgoing O2 angle theta, in terms of hbar, c, theta, Delta U, and m. Reusable conclusions: For theta <= pi/2, omega\_min = 3*m*c\textasciicircum{}2*[1 - sqrt(1 - (Delta U/(3*m*c\textasciicircum{}2))*(2*sin(theta)\textasciicircum{}2 + 1))]/[hbar*(2*sin(theta)\textasciicircum{}2 + 1)]. For theta >= pi/2 use the same threshold evaluated at theta = pi/2. Policy: natural\_language\_prerequisite\_only; do\_not\_import\_Lean\_output
\end{itemize}

\paragraph{Formalization target.}
create a compiling Lean file with sorry bodies at `IPhO2026Problems/problem\_IPhO\_2026\_1\_C\_2.lean`.
The Lean declarations must preserve the physical quantities, dimensions or dimensional roles, figure labels, governing-law hypotheses, and final relation expressed by this problem.
Use Mathlib/Physlib names found through LeanExplore where available. If a domain API is missing, introduce faithful local abstractions rather than scalar placeholder aliases.

\begin{theorem}[Physics formalization target]
\label{thm:physics:IPhO_2026_1_C_2:target}
The assigned autoformalize agent should translate this physics problem into Lean declarations in the covered file, with theorem and lemma proof bodies written as `by sorry`.
\end{theorem}
\begin{proof}
This is an autoformalization task, not a proof task. Produce faithful statements that can later be proved without weakening the source contract.
\end{proof}
% --- Archon physics formalization source end ---
